\section{Dataset Description, Cleaning, and Exploratory Analysis}

The project uses the HAM10000 dermoscopic image collection \citep{tschandl2018ham10000}. The active pipeline reads the original image archive and \texttt{HAM10000\_metadata.csv}, reorganizes images into class folders, and loads them through \texttt{ImageFolder}. The alternative low-resolution CSV representations stored in the archive are not used by the current code.

\begin{table}[H]
    \centering
    \caption{Dataset summary from the current repository.}
    \label{tab:datasetsummary}
    \begin{tabular}{>{\raggedright\arraybackslash}p{0.30\textwidth} >{\raggedright\arraybackslash}p{0.60\textwidth}}
        \toprule
        Item & Current dataset evidence \\
        \midrule
        Source dataset & HAM10000 \\
        Original inputs & Two image folders and one metadata CSV \\
        Active label field & \texttt{dx} \\
        Final organization & \texttt{archive/Training} and \texttt{archive/Testing} class folders \\
        Split method & Grouped by \texttt{lesion\_id}, stratified by \texttt{dx} \\
        Training images & 8{,}011 \\
        Testing images & 2{,}004 \\
        Training lesions & 5{,}975 unique lesions \\
        Testing lesions & 1{,}495 unique lesions \\
        Leakage check & 0 image overlap, 0 lesion overlap \\
        Original image size & \(600\times450\) \\
        Metadata used & \texttt{lesion\_id}, \texttt{image\_id}, \texttt{dx}, \texttt{dx\_type}, \texttt{age}, \texttt{sex}, \texttt{localization} \\
        \bottomrule
    \end{tabular}
\end{table}

The preparation script performs practical cleaning and reorganization rather than deep manual curation. It validates the input paths, reads the metadata file, builds a grouped split, recreates the output directories, and copies images into diagnosis-specific folders. Based on the available code and regenerated statistics, the current pipeline removes lesion-level leakage and aligns folder names with metadata labels. No explicit evidence of corrupted-image filtering, manual relabeling, or duplicate removal beyond grouped splitting was found in the repository.

\begin{table}[H]
    \centering
    \caption{Seven lesion categories used in the classifier.}
    \label{tab:categories}
    \begin{tabular}{ll}
        \toprule
        Label & Full class name \\
        \midrule
        akiec & Actinic keratosis / intraepithelial carcinoma \\
        bcc   & Basal cell carcinoma \\
        bkl   & Benign keratosis-like lesion \\
        df    & Dermatofibroma \\
        mel   & Melanoma \\
        nv    & Melanocytic nevus \\
        vasc  & Vascular lesion \\
        \bottomrule
    \end{tabular}
\end{table}

The regenerated EDA shows strong class imbalance, with the nevus category dominating both training and testing, while Figure~\ref{fig:splitintegrity} confirms that train and test remain separate at both image and lesion levels. These properties directly affect evaluation because they reduce leakage but still leave a difficult class-frequency imbalance.

\begin{figure}[H]
    \centering
    \includegraphics[width=0.84\textwidth]{class_distribution.png}
    \caption{Class distribution in the grouped split.}
    \label{fig:classdist}
\end{figure}

\begin{figure}[H]
    \centering
    \includegraphics[width=0.80\textwidth]{split_integrity.png}
    \caption{Grouped split integrity check.}
    \label{fig:splitintegrity}
\end{figure}

Figure~\ref{fig:gallery} shows representative examples from the seven categories. Several classes overlap in color, border quality, and internal texture, which helps explain why per-class evaluation remains necessary even when overall accuracy is high.

\begin{figure}[H]
    \centering
    \includegraphics[width=0.70\textwidth]{sample_gallery.png}
    \caption{Representative training examples from the seven classes.}
    \label{fig:gallery}
\end{figure}
